\section{Emerging Applications}
\label{app:applications}
 
We highlight three domains where verbal RL has shown the strongest impact. In each, we note which pillars are most active.
 
\subsection{Coding Assistants}
 
Coding assistants exemplify the interplay of all three pillars. Grounding signal (Pillar~3) enters through natural-language issue descriptions and test suites that define the task. Deliberative feedback (Pillar~1) drives the iterative loop: agents translate goals into code, execute it, and refine outputs using compiler errors, stack traces, test failures, and linter warnings as verbal feedback \citep{yang2024swe, chen2024teaching}. Modern systems autonomously handle debugging, test generation, repository navigation, and multi-step issue resolution \citep{yang2023intercode}. Tool outputs such as search and dependency traces enable reasoning across large codebases, while persistent memory captures user preferences and repository constraints \citep{wang2023voyager}. When successful and failed trajectories are collected for fine-tuning, they become a learning signal (Pillar~2). SWE-bench \citep{jimenez2024swebenchlanguagemodelsresolve} has emerged as the standard benchmark for evaluating these agents.
 
\subsection{Scientific Discovery and Research Workflows}
 
Scientific discovery primarily leverages Pillar~1, since the scientific process itself, hypothesizing, experimenting, interpreting, and revising, is inherently language-driven and iterative. In drug discovery, agents refine molecular candidates using verbal feedback from chemistry and simulation tools \citep{m2024augmenting}. In materials science, agents explore compositional spaces through simulator-guided refinement \citep{szymanski2023autonomous}. Beyond individual tasks, automated systems now support end-to-end workflows including ideation, literature review, experiment design, execution, and manuscript generation \citep{lu2024ai}. Human-in-the-loop frameworks further incorporate researcher feedback during iterative refinement \citep{schmidgall2025agent}. When tool outputs define experimental constraints or success criteria, Pillar~3 (grounding) is also active.
 
\subsection{Robotics}
 
Robotics is the canonical application of Pillar~3, where language must ground in physical action. SayCan \citep{ahn2022can} translates high-level instructions into grounded robot actions constrained by physical capabilities, while Inner Monologue \citep{huang2022inner} incorporates environmental feedback for replanning and recovery during execution (Pillar~1). Code as Policies \citep{liang2022code} converts language instructions into executable control code and improves behavior through execution feedback. More recent systems such as ELLMER \citep{mon2025embodied} demonstrate long-horizon task completion in dynamic environments through continual adaptation to verbal and environmental feedback. Industry systems are beginning to deploy these capabilities: NVIDIA's GR00T N1 \citep{bjorck2025gr00t} provides a foundation model for language-conditioned robot control, while 1X Technologies' NEO humanoid learns tasks through verbal interaction and human guidance.\footnote{\url{https://www.1x.tech}} Together, these systems reflect a broader shift toward robotic agents that ground language in open-ended physical environments.


% \section{Emerging Applications}
% \label{app:applications}

% We highlight six domains where verbal RL has been most visible. For each domain, we identify the pillars that are most active.

% \subsection{Coding Assistants}

% Coding assistants bring together all three pillars. Grounding signal (Pillar~3) enters through natural-language issue descriptions, repository context, and test suites that define the task. Deliberative feedback (Pillar~1) drives the iteration loop: agents translate goals into code, execute their proposed changes, and revise them using compiler errors, stack traces, test failures, and linter warnings as textual feedback \citep{yang2024swe, chen2024teaching}. Recent systems have moved beyond single-turn code generation toward repository-level agents that support debugging, test generation, multi-file navigation, and multi-step issue resolution \citep{yang2023intercode}. Systems such as Claude Code, OpenAI Codex, and GitHub Copilot coding agents now operate in development harnesses where repository state, execution feedback, and collaboration workflows shape the agent's next actions \citep{ning2026code}. Tool outputs, including search results, dependency traces, and CI/CD logs, help agents reason across large codebases, while persistent memory can record user preferences and repository-specific constraints \citep{wang2023voyager}. When successful and failed trajectories are collected for fine-tuning, they also become a learning signal (Pillar~2). In this setting, real development traces are increasingly useful as training data for improving later coding models \citep{ning2026code}. SWE-bench \citep{jimenez2024swebenchlanguagemodelsresolve} has become a standard benchmark, while newer benchmarks such as SWE-Lancer and Terminal-Bench extend evaluation to freelance-style tasks and interactive terminal environments.

% \subsection{Scientific Discovery and Research Workflows}

% Scientific discovery relies heavily on Pillar~1 because scientific work is naturally iterative: researchers form hypotheses, run experiments, interpret results, and revise their assumptions. In drug discovery, agents refine molecular candidates using feedback from chemistry and simulation tools \citep{m2024augmenting}. In materials science, agents explore compositional spaces through simulator-guided refinement \citep{szymanski2023autonomous}, and self-driving laboratories such as Berkeley's A-Lab have shown how synthesis, measurement, and revision can be connected in a single experimental loop. In formal and computational settings, systems such as AlphaProof connect language models with executable verifiers, so that candidate proof steps can be checked before being reused as feedback \citep{ning2026code}. Beyond individual tasks, automated systems now assist with broader research workflows, including ideation, literature review, experiment design, execution, and manuscript drafting \citep{lu2024ai}. Human-in-the-loop frameworks add another feedback source by allowing researchers to guide or correct the system during revision \citep{schmidgall2025agent}. When tool outputs define experimental constraints or success criteria, Pillar~3 is also active. Benchmarks such as MLAgentBench, ScienceAgentBench, and DiscoveryBench are beginning to provide standardized evaluation for scientific agents across domains \citep{ning2026code}.

% \subsection{Robotics}

% Robotics is a central application of Pillar~3, since language instructions must be grounded in physical actions. SayCan \citep{ahn2022can} maps high-level instructions to robot actions while accounting for physical feasibility, and Inner Monologue \citep{huang2022inner} uses environmental feedback for replanning and recovery during execution (Pillar~1). Code as Policies \citep{liang2022code} converts language instructions into executable control code and improves behavior through execution feedback. A distinctive feature of embodied agents is that code can serve two purposes: it grounds actions in executable procedures, and it provides a reusable representation for skills that can persist across tasks. Voyager \citep{wang2023voyager} introduced this idea through a growing skill library in Minecraft, and later work has extended similar patterns to tabletop manipulation, human-corrected plans, and continual learning \citep{ning2026code}. More recent systems such as ELLMER \citep{mon2025embodied} study long-horizon task completion in dynamic environments, where agents must adapt to verbal and environmental feedback over time. Industry systems are also exploring this direction: NVIDIA's GR00T N1 \citep{bjorck2025groot} provides a foundation model for language-conditioned robot control, while 1X Technologies' NEO humanoid is designed to learn tasks through verbal interaction and human guidance.\footnote{\url{https://www.1x.tech}}

% \subsection{GUI/OS Automation}

% GUI and operating system automation is a setting where all three pillars can appear together. The environment is built from executable software: observations are rendered from code, such as HTML, CSS, and accessibility APIs, and actions are issued through code-level interfaces, such as DOM events, shell commands, and keystrokes \citep{ning2026code}. Grounding signal (Pillar~3) is especially important because agents must interpret screen states from DOM trees, accessibility labels, or vision-language descriptions and map them to executable actions. Deliberative feedback (Pillar~1) drives the interaction loop: an agent attempts an action, observes the resulting page state, error dialog, or navigation failure, and revises its plan. When interaction traces are collected to train better navigation policies, they also provide a learning signal (Pillar~2). Benchmarks such as WebArena, OSWorld, and WindowsAgentArena provide standardized evaluation in sandboxed web, desktop, and mobile environments. Systems such as Anthropic's Computer Use, OpenAI's Operator, and Google DeepMind's Project Mariner show that similar perception-action-feedback patterns can be applied to browser and desktop control outside benchmark settings \citep{ning2026code}. A central challenge in this domain is choosing the right level of granularity: a user instruction may require a single click or a long workflow, and the feedback from the interface must guide the agent at the appropriate level of abstraction.

% \subsection{DevOps and CI/CD Workflows}

% DevOps pipelines extend the coding-assistant setting from source-code edits to infrastructure-level feedback. Build logs, deployment errors, monitoring alerts, security scan results, and pull-request review comments all provide textual signals that agents can use to diagnose and repair failures. Deliberative feedback (Pillar~1) is primary: agents receive execution traces from CI/CD systems, interpret likely failure modes, and propose fixes using evidence from the surrounding infrastructure. Grounding signal (Pillar~3) is also present because deployment specifications, infrastructure-as-code templates, and service-level objectives define the task constraints in natural language or structured text. Recent systems demonstrate this pattern at production scale. For example, LingmaAgent at Alibaba Cloud resolves a substantial share of in-house issues through feedback from build and deployment systems \citep{ning2026code}. Multi-agent architectures further divide responsibilities across planner, coder, tester, and reviewer roles operating over shared repositories, with the CI pipeline serving as both a feedback source and a stopping condition. The main verbal RL challenge in this domain is feedback disambiguation: infrastructure failures, network errors, and genuine code bugs can produce similar error messages, but they require different recovery strategies.

% \subsection{Data Analysis and Visualization}

% Data analysis offers a compact verbal RL loop: agents generate analytical code, execute it, inspect the result, and revise their approach. The feedback is often textual or easily rendered as text, including statistical summaries, table previews, error messages, and chart descriptions (Pillar~1). Grounding signal (Pillar~3) enters through natural-language task descriptions that specify the analytical objective, data constraints, and required output format. Code-interpreter-style agents illustrate this pattern: they write Python or SQL, run it in a sandboxed environment, observe the output, and iteratively refine the analysis. Benchmarks such as MLAgentBench \citep{ning2026code} evaluate agents on open-ended machine-learning research tasks that require reading code, running experiments, and improving metrics over successive iterations. The distinguishing feature of this domain is the short feedback cycle. Unlike scientific discovery, where experiments may take hours or days, many data-analysis tasks provide execution feedback within seconds, allowing multiple revisions in a single session. The verbal RL challenge is output interpretation: agents must decide whether a statistical result, visualization, or model metric reflects genuine progress toward the analytical goal or is merely an artifact of the data split, preprocessing choice, or hyperparameter setting.


% \section{Subcategory Summary Tables}
% \label{sec:subcat-tables}


% \begin{table*}[!th] % The * makes it span both columns
%     \centering
%     \small
%     \renewcommand{\arraystretch}{1.3}
%     \caption{Pillar 1---Language as Deliberative Feedback.}
%     \label{tab:pillar1-subcat}
%     \begin{tabularx}{\linewidth}{@{} >{\raggedright\arraybackslash}p{0.22\linewidth} >{\raggedright\arraybackslash}X >{\raggedright\arraybackslash}p{0.35\linewidth} @{}}
%         \toprule
%         \textbf{Subcategory} & \textbf{Description} & \textbf{Representative Papers} \\
%         \midrule
%         Self-Refinement / Critique & Model iteratively revises its output using its own feedback & \citet{madaan2023selfrefine}; \citet{shinn2023reflexion}; \citet{welleck2023generating} \\
%         \midrule
%         \addlinespace
%         Externally Grounded Critique & Critique grounded in tools, execution, retrieval, or environment & \citet{gou2024critic}; \citet{yao2023react} \\
%         \midrule
%         \addlinespace
%         Multi-Agent Debate & Multiple agents challenge each other's outputs & \citet{du2023debate}; \citet{bai2022constitutional} \\
%         \midrule
%         \addlinespace
%         Reflective Memory & Feedback from prior attempts stored as verbal lessons & \citet{shinn2023reflexion}; \citet{park2023genagents} \\
%         \midrule
%         \addlinespace
%         Search-Guided Deliberation & Explore multiple reasoning paths with language-guided search & \citet{yao2023tree}; \citet{hao2023reasoning} \\
%         \bottomrule
%     \end{tabularx}
% \end{table*}



% \begin{table*}[!th]
%     \centering
%     \small
%     \renewcommand{\arraystretch}{1.3}
%     \caption{Pillar 2---Language as Learning Signal.}
%     \label{tab:pillar2-subcat}
%     \begin{tabularx}{\linewidth}{@{} >{\raggedright\arraybackslash}p{0.22\linewidth} >{\raggedright\arraybackslash}X >{\raggedright\arraybackslash}p{0.35\linewidth} @{}}
%         \toprule
%         \textbf{Subcategory} & \textbf{Description} & \textbf{Representative Papers} \\
%         \midrule
%         Self-Improvement & Model bootstraps better training data from own critiques & \citet{zelikman2022star}; \citet{huang2022selfimprove} \\
%         \midrule
%         \addlinespace
%         Preference Shaping & NL feedback $\rightarrow$ preferences $\rightarrow$ reward-model signals & \citet{ouyang2022training}; \citet{rafailov2023direct}; \citet{lee2023rlaif} \\
%         \midrule
%         \addlinespace
%         Feedback-Conditioned Modeling & Model trained to condition on feedback/corrections directly & \citet{scheurer2023training}; \citet{luo2025fcp} \\
%         \bottomrule
%     \end{tabularx}
% \end{table*}

% % ==========


% \begin{table*}[!th]
%     \centering
%     \small
%     \renewcommand{\arraystretch}{1.3}
%     \caption{Pillar 3---Language as Grounding Signal.}
%     \label{tab:pillar3-subcat}
%     \begin{tabularx}{\linewidth}{@{} >{\raggedright\arraybackslash}p{0.22\linewidth} >{\raggedright\arraybackslash}X >{\raggedright\arraybackslash}p{0.35\linewidth} @{}}
%         \toprule
%         \textbf{Subcategory} & \textbf{Description} & \textbf{Representative Papers} \\
%         \midrule
%         Goal Grounding & Language specifies what the agent should optimize & \citet{ahn2022can}; \citet{chevalier2018babyai}; \citet{andreas2017modular} \\
%         \midrule
%         \addlinespace
%         State Grounding & Language describes agent's current situation & \citet{yao2020keep}; \citet{shridhar2020alfworld}; \citet{hausknecht2020interactive} \\
%         \midrule
%         \addlinespace
%         Action Grounding & Language constrains or guides the action space & \citet{huang2022inner}; \citet{sharma2022correcting}; \citet{lynch2021language} \\
%         \midrule
%         \addlinespace
%         Reward Code Generation & LLM writes executable reward function from NL specification & \citet{ma2024eureka}; \citet{xie2024text2reward}; \citet{yu2023language} \\
%         \bottomrule
%     \end{tabularx}
% \end{table*}